\documentclass[11pt,letterpaper]{article}

\usepackage[T1]{fontenc}
\usepackage{amsmath,amssymb}
\usepackage{tikz}
\usepackage{pgfplots}
\pgfplotsset{compat=1.18}
\usepackage{algorithm}
\usepackage{algpseudocode}
\usepackage{booktabs}
\usepackage{enumitem}
\usepackage[hidelinks]{hyperref}

\newcommand{\N}{\mathbb{N}}
\newcommand{\Z}{\mathbb{Z}}
\newcommand{\Q}{\mathbb{Q}}
\newcommand{\R}{\mathbb{R}}
\newcommand{\C}{\mathbb{C}}
\newcommand{\Map}[3]{#1 \colon #2 \to #3}
\newcommand{\Mod}[1]{#1\text{-}\mathrm{Mod}}
\newcommand{\GroupHom}{\mathrm{Hom}_{\mathbf{Grp}}}
\newcommand{\RingHom}{\mathrm{Hom}_{\mathbf{Ring}}}

\newtheorem{theorem}{Theorem}
\newtheorem{lemma}[theorem]{Lemma}
\newtheorem{corollary}[theorem]{Corollary}
\newtheorem{definition}[theorem]{Definition}
\newtheorem{remark}[theorem]{Remark}
\newenvironment{proof}{\noindent\textit{Proof.}}{\hfill\(\square\)}

\newcommand{\OrderNorm}{\operatorname{N}}
\newcommand{\RawMap}{T}
\newcommand{\OddPart}{\operatorname{odd}}

\title{Exact Division, Unit Phases, and Escape in Quadratic Orders}
\author{}
\date{\today}

\begin{document}
\maketitle

\begin{abstract}
Exact division in an imaginary quadratic order is decided by two integer
congruences. A unimodular triangular basis gives precisely the norm of the
divisor many residue classes. These certificates retain the division guard,
including in orders without elementwise unique factorization.
For the family \(q=2\), \(c=0\), \(d=1+\sqrt{-D}\), every nonzero trajectory
is eventually periodic with primitive period twelve when \(D=1\) or \(D=3\);
for every other positive integer \(D\), every nonzero trajectory escapes
along the full raw sequence. First-division densities and an odd-offset
dyadic family are also determined explicitly.
\end{abstract}

\section{The order and the exact branch}

The coordinate lemmas have Lean proofs in \cite{LibraryCoordinates26};
the dyadic guard and restricted block lemmas are in \cite{LibraryDyadic26}.
The finite arithmetic checks are implemented in \cite{LibraryQuadraticChecks26}.
The trajectory classifications proved below are not yet formalized in those Lean files.

Let \(\tau,\nu\in\Z\), assume \(\tau^2<4\nu\), and set
\[
w=\frac{\tau+i\sqrt{4\nu-\tau^2}}2,\qquad
R=\Z[w]\subset\C.
\]
The coordinates \(x+yw\), \(x,y\in\Z\), are unique. The relation is
\(w^2=\tau w-\nu\). Conjugation gives
\[
\OrderNorm(x+yw)=|x+yw|^2=x^2+\tau xy+\nu y^2.
\]
For \(q,c,d\in R\), with \(d\ne0\), define
\[
\RawMap_{q,c,d}(z)=
\begin{cases}
z/d,&d\mid z\text{ in }R,\\
qz+c,&d\nmid z\text{ in }R.
\end{cases}
\]
Division denotes the unique exact quotient in \(R\). Its test is not
replaced by divisibility in \(\C\). The zero state is fixed.
Replacing \(d\) by a unit associate preserves its ideal but can change
the quotient and the raw trajectory. Enlarging the order can change even
the guard: in \(\Z[\sqrt{-3}]\), \(1+\sqrt{-3}\nmid2\), whereas in the
Eisenstein order, with \(\omega=(-1+\sqrt{-3})/2\), its quotient is
\((1-\sqrt{-3})/2=-\omega\). Thus the larger-order map need not restrict
to the original raw map.

\begin{lemma}[Positive and multiplicative norm]
For \(z\ne0\), \(\OrderNorm(z)\) is a positive integer. For all \(z,u\in R\),
\(\OrderNorm(zu)=\OrderNorm(z)\OrderNorm(u)\).
\end{lemma}
\begin{proof}
The identity
\[
4\OrderNorm(x+yw)=(2x+\tau y)^2+(4\nu-\tau^2)y^2
\]
is positive when \(y\ne0\); when \(y=0\) and \(x\ne0\), the norm is \(x^2>0\).
The coefficient product is
\[
(x+yw)(s+tw)=(xs-\nu yt)+(xt+ys+\tau yt)w.
\]
Substitution in the norm form and expansion proves multiplicativity.
Equivalently, conjugation multiplies and
\((zu)\overline{zu}=(z\bar z)(u\bar u)\).
\end{proof}

\section{Integral division and a finite quotient}

\begin{theorem}[Adjugate division certificate]
Let \(d=s+tw\ne0\) and \(n=\OrderNorm(d)\). Then
\[
d\mid x+yw
\iff
n\mid (s+\tau t)x+\nu ty
\quad\text{and}\quad
n\mid -tx+sy.
\]
When these conditions hold, the exact quotient is
\[
\frac{x+yw}{d}
=\frac{(s+\tau t)x+\nu ty}{n}
+\frac{-tx+sy}{n}\,w.
\]
\end{theorem}
\begin{proof}
Multiplication by \(d\) has matrix
\[
M_d=\begin{pmatrix}s&-\nu t\\t&s+\tau t\end{pmatrix},
\qquad
\det M_d=n>0.
\]
Its adjugate is
\(\left(\begin{smallmatrix}s+\tau t&\nu t\\-t&s\end{smallmatrix}\right)\).
The rational solution of \(M_d(u,v)^{\mathsf T}=(x,y)^{\mathsf T}\) is
unique and has the displayed coordinates. It is integral precisely
under the two stated congruences. Multiplying the adjugate by \(M_d\)
also proves uniqueness among integral quotients.
\end{proof}

\begin{theorem}[Triangular representatives]
Let \(g=\gcd(s,t)>0\), and choose \(u,v\in\Z\) with \(su+tv=g\).
Put
\[
L=n/g,\qquad h=sv-(\tau s+\nu t)u.
\]
Then \(dR\) has additive basis \((L,0),(h,g)\). Each residue class has
exactly one representative \(x_0+y_0w\) with
\(0\le x_0<L\) and \(0\le y_0<g\). Consequently,
\[
|R/dR|=Lg=n.
\]
\end{theorem}
\begin{proof}
The Euclidean algorithm on \(|s|,|t|\) terminates because each nonzero
remainder is smaller than the preceding positive divisor. Its
coefficient recurrence, with the signs of \(s,t\) restored, supplies
the stated identity. Since \(g\mid s,t\), the matrix
\[
U=\begin{pmatrix}(s+\tau t)/g&v-\tau u\\-t/g&u\end{pmatrix}
\]
is integral and has determinant \((su+tv)/g=1\). Direct multiplication gives
\[
M_dU=\begin{pmatrix}L&h\\0&g\end{pmatrix}.
\]
The integral inverse of \(U\) proves equality of the generated lattices.
For an arbitrary point, set
\[
y_0=y\bmod g,\qquad k=(y-y_0)/g,\qquad
x_0=(x-kh)\bmod L.
\]
Its difference from \((x_0,y_0)\) is an integral combination of the two
basis vectors. If two canonical representatives differ by such a
combination, their vertical difference is a multiple of \(g\) of
absolute value less than \(g\), hence zero. Their horizontal difference
is then a multiple of \(L\) of absolute value less than \(L\), hence zero.
This proves both exhaustion and uniqueness, including \(s=0\), \(t=0\),
and the one-element quotient when \(d\) is a unit.
\end{proof}

\section{The guard in a nonfactorial order}

In \(R_5=\Z[w]\), \(w^2=-5\), set \(q=2\), \(d=1+w\).
The equality \((1+w)(1-w)=6\) gives \(d\mid6\), whereas
\(3/d=(1-w)/2\notin R_5\). Thus
\[
d\mid q\cdot3,\qquad d\nmid3.
\]
Every common divisor \(a\) of \(2\) and \(d\) has norm dividing both
\(4\) and \(6\). A nonunit would therefore have norm two, but
\(x^2+5y^2=2\) has no integral solution. Hence absence of a common
nonunit element divisor does not justify guard cancellation.

The map \(x+yw\mapsto x+y\bmod2\) is a ring homomorphism to \(\Z/2\Z\).
It sends \(2,d\) to zero and \(1\) to one, proving that no identity
\(a\cdot2+bd=1\) exists. By contrast, a supplied identity
\(aq+bd=1\) proves
\[
d\mid qz\ \Longrightarrow\
d\mid a(qz)+b(dz)=z.
\]
This is the exact sufficient cancellation hypothesis.

\section{A finite certificate for division-free escape}

\begin{theorem}[Affine residue certificate]
Assume \(\OrderNorm(q),\OrderNorm(d)\ge2\).
There exists a division-free escaping raw trajectory if and only if
\[
F:R/dR\longrightarrow R/dR,\qquad [z]\longmapsto[qz+c]
\]
has a directed cycle avoiding the zero class.
If it has no such cycle, every raw seed reaches a divisible state at
some index \(j<\OrderNorm(d)\).
\end{theorem}
\begin{proof}
The affine map descends because its difference on \(z+d\ell,z\) is
\(qd\ell\). A division-free trajectory supplies an infinite path among
the finitely many nonzero classes, hence a repeated vertex and a cycle.
For recovery from a cycle, lift a vertex to \(r+d\) and \(r+2d\).
Both affine paths retain the same zero-avoiding residue itinerary.
Writing \(\delta(z)=(q-1)z+c\), their displacements differ by
\((q-1)d\ne0\). At least one displacement has modulus at least
\(|q-1||d|/2\). For that lift,
\[
\RawMap^j(z)=q^j\frac{\delta(z)}{q-1}-\frac c{q-1},\qquad
|\RawMap^j(z)|\ge\frac{|d|}{2}|q|^j-\frac{|c|}{|q-1|}
\longrightarrow\infty.
\]
Every step is affine because the same residue cycle avoids zero.
Finally, if the first \(\OrderNorm(d)\) states were nondivisible,
two of their classes would coincide among only
\(\OrderNorm(d)-1\) nonzero classes. Their intervening affine path
would give a zero-avoiding cycle. Before the first division, the raw
and affine paths are identical.
\end{proof}

\section{Finite division budgets}

For a prime \(p\) and a positive integer \(n\), write
\(v_p(n)=\max\{a\in\N:p^a\mid n\}\).

\begin{theorem}[Norm-prime budget]
Assume \(c=0\), \(\OrderNorm(q),\OrderNorm(d)\ge2\). Let \(p\) be a rational
prime with \(p\mid\OrderNorm(d)\) and \(p\nmid\OrderNorm(q)\).
For a nonzero seed \(z\), set
\[
e=v_p(\OrderNorm(d)),\quad
V=\left\lfloor\frac{v_p(\OrderNorm(z))}{e}\right\rfloor,\quad
M=\OrderNorm(d).
\]
There are at most \(V\) division instructions. No division occurs at
a time index at least \(K_0=M(V+1)\), and for \(j\ge K_0\),
\[
\OrderNorm(\RawMap^j(z))\ge\OrderNorm(q)^{j-K_0}\ge2^{j-K_0}.
\]
\end{theorem}
\begin{proof}
No nonzero state becomes zero. At an affine instruction the norm
valuation is unchanged; at a division it decreases by exactly \(e\).
Thus it is \(v_p(\OrderNorm(z))-eB_j\), where \(B_j\) counts divisions
before time \(j\), and nonnegativity gives \(B_j\le V\).
From any state, either its next division is executed within \(M\)
ticks, or the first \(M\) states are nondivisible. In the latter case
their repeated affine residue supplies an indefinitely zero-avoiding
path, so no later division occurs. Successive divisions therefore
consume at most \(M\) ticks each. After at most \(V\) such blocks,
one further block proves that the tail is entirely affine.
The norm at time \(K_0\) is a positive integer; norm multiplicativity
then proves the bound. The assumption \(c=0\) is essential to the
valuation identity.
\end{proof}

\section{\texorpdfstring{The uniform \(D\)-family}{The uniform D-family}}

Set \(D\ge1\), \(w=\sqrt{-D}\), \(R_D=\Z[w]\),
\(q=2\), \(d=1+w\). Put \(M=D+1=2^e m\), with \(m\) odd.
The quotient is explicitly
\[
R_D/(1+w)\cong\Z/M\Z,\qquad [x+yw]\longmapsto[x-y].
\]
Indeed, \(dR_D\) has basis \((M,0),(1,1)\), and substitution \(w=-1\)
respects \(w^2=-D\) modulo \(M\).

\begin{theorem}[First division and its density]
For \(c=0\) and an initial integer \(n\ne0\), a first division exists
if and only if \(m\mid n\). If it exists, its state index is
\[
a=\max\{0,e-v_2(|n|)\}.
\]
The positive-integer density of seeds reaching a first division is
\(1/m\). Every other positive seed follows \(2^j n\) and escapes.
\end{theorem}
\begin{proof}
Before the first division the states are \(2^j n\).
The guard is \(M\mid2^j n\), equivalent to \(m\mid n\) and
\(e\le j+v_2(|n|)\). This gives the existence and least index.
The number of multiples of \(m\) among \(1,\ldots,X\) is
\(\lfloor X/m\rfloor\); division by \(X\) proves the density.
The zero seed divides at index zero and is fixed.
\end{proof}

\begin{lemma}[Positive-power dependence]
The equation
\[
2^a=(1+w)^b,\qquad a,b\ge1,
\]
has an integral solution if and only if \(D\in\{1,3\}\).
\end{lemma}
\begin{proof}
Suppose equality holds. Conjugating and dividing gives
\(u^b=1\), where \(u=(1+w)/(1-w)\), and
\[
t=u+u^{-1}=\frac{2(1-D)}{1+D}\in\Q.
\]
Define integral polynomials \(P_0=2,P_1=X\) and
\(P_{j+1}=XP_j-P_{j-1}\). Induction gives
\(P_j(t)=u^j+u^{-j}\), and \(P_j\) is monic of degree \(j\) for \(j\ge1\).
Thus \(P_b(t)-2=0\) is a monic integral equation.
A rational root of a monic integral polynomial is integral: writing
it in lowest terms and clearing denominators forces its denominator
to divide its numerator. Hence \(t\in\Z\).
For \(D\ge1\), \(-2<t\le0\), so \(t=-1\) or \(t=0\).
These equations give \(D=3\) or \(D=1\), respectively.
Conversely, when \(D=1\), \(d^8=16=2^4\); when \(D=3\),
\(d^6=64=2^6\), by expansion using \(w^2=-D\).
\end{proof}

\begin{theorem}[Dyadic counts and escape exponent]
Let \(D=2^e-1\), \(e\ge3\), \(q=2,c=0,d=1+w\), and let
\(z\ne0\). Put \(V=v_2(\OrderNorm(z))\). If \(B_j\) counts divisions
in the first \(j\) raw ticks, then
\[
B_j\le\left\lfloor\frac{j+V}{e}\right\rfloor,
\qquad
\OrderNorm(\RawMap^j(z))\ge
\OrderNorm(z)\,2^{((e-2)j-(e+2)V)/e}.
\]
After at most \(L=5V+7e+1\) ticks, the instruction word is periodic
with primitive period \(e\), containing \(e-1\) affine instructions
and one division. The arithmetic states still escape. In particular,
\[
\lim_{j\to\infty}\frac{B_j}{j}=\frac1e,
\qquad
\lim_{j\to\infty}\frac{\log|\RawMap^j(z)|}{j}
=\frac{e-2}{2e}\log2>0.
\]
\end{theorem}
\begin{proof}
Construct a square root in \(\Z_2=\varprojlim_k\Z/2^k\Z\) rather than
change the original order or its guard. Set
\[
u_0=1,\qquad u_{n+1}=1+2^{e-2}u_n^2.
\]
Every \(u_n\) is odd. The first difference is divisible by
\(2^{e-2}\); subsequent differences satisfy
\[
u_{n+1}-u_n=2^{e-2}(u_n-u_{n-1})(u_n+u_{n-1}).
\]
Since the last factor is even, the stabilized precision increases
by at least \(e-1\) bits per iteration. The residues therefore define
a unique limit \(u\in\Z_2\), with
\(u=1+2^{e-2}u^2\) and \(u\) odd. For
\[
r=-1+2^{e-1}u
\]
we obtain
\[
r^2=1-2^e=-D,\quad v_2(1+r)=e-1,\quad v_2(1-r)=1.
\]
Here a nonzero \(2\)-adic integer has a finite valuation: factor out
its first nonzero binary digit. Odd factors multiply to odd factors,
which proves additivity of valuation under nonzero products.

The two ring homomorphisms
\[
\Phi_+(x+yw)=x+yr,\qquad \Phi_-(x+yw)=x-yr
\]
send \(d\) to elements of valuations \(e-1\) and one. Neither sends
a nonzero order element to zero: their product is the positive
integer \(x^2+Dy^2\), and integer embedding in \(\Z_2\) is injective.
Let \(h_j,\ell_j\) be the respective valuations of the actual raw
state. Their sum initially is \(V\), so \(h_0,\ell_0\le V\).
Each legal affine step adds one to both valuations. Each legal
division subtracts \(e-1\) from \(h_j\) and one from \(\ell_j\).
Consequently,
\[
h_j=h_0+j-eB_j\ge0,\qquad
\ell_j=\ell_0+j-2B_j\ge0.
\]
The first identity proves the count bound. The norm identity
\[
\OrderNorm(\RawMap^j(z))=\OrderNorm(z)2^{2j-(e+2)B_j}
\]
then proves the quantitative growth bound. These computations use the
same global instruction at both homomorphisms.

The valuation pair also determines that instruction exactly:
\[
\begin{gathered}
\text{division}\quad\Longleftrightarrow\\
\bigl(h_j=e-1\text{ and }\ell_j=1\bigr)
\quad\text{or}\quad
\bigl(h_j\ge e\text{ and }\ell_j\ge2\bigr).
\end{gathered}
\]
To prove this, the rational quotient has coordinates
\((x+Dy)/2^e\) and \((y-x)/2^e\). They are dyadic rationals.
Its two images are \(\Phi_+(z)/(1+r)\) and
\(\Phi_-(z)/(1-r)\). Integral coordinates require these images
to be integral with equal parity. Conversely, integral images of
equal parity give integral \(2\)-adic coordinates by taking their
sum divided by two and their difference divided by \(2r\);
\(r\) is odd. A dyadic rational integral at two is an integer, so
this condition recovers the exact quotient in \(R_D\).
The two images have equal parity precisely when they are both odd
or both even, giving the displayed criterion.

Set \(K=2V+3e\). The count bound gives, for every \(j\ge K\),
\[
\ell_j\ge\frac{(e-2)j-2V}{e}\ge3.
\]
The division criterion is now simply \(h_j\ge e\). On this tail,
\(h_j\) increases by one below \(e\) and decreases by \(e-1\)
at or above \(e\). Since \(h_K\le V+K\), finitely many consecutive
decreases, and at most one increase from zero, put it in the cycle
\[
1\longrightarrow2\longrightarrow\cdots\longrightarrow e
\longrightarrow1.
\]
This occurs by \(2K+V+1\le L\). It has exactly one division per
\(e\) ticks, and that single distinguished instruction proves the
primitive word period \(e\). Hence \(B_j/j\to1/e\).
Substitution in the norm identity proves the stated logarithmic limit.
Indeed, every tail block of \(e\) ticks satisfies the stronger identity
\[
\RawMap^{j+e}(z)=\frac{2^{e-1}}d\RawMap^j(z),\qquad
\OrderNorm(\RawMap^{j+e}(z))=2^{e-2}\OrderNorm(\RawMap^j(z)).
\]
Although the scalar \(2^{e-1}/d=(1-w)/2\) need not belong to the order,
the displayed outputs do: they are obtained by the same legal raw word.
\end{proof}

\begin{corollary}[Odd-norm seeds]
Under the preceding dyadic hypotheses, if \(\OrderNorm(z)\) is odd,
then \(B_0=0\) and
\[
B_j=\left\lfloor\frac{j-1}{e}\right\rfloor\qquad(j\ge1).
\]
\end{corollary}
\begin{proof}
Initially \(h_0=\ell_0=0\). The first tick is affine.
The next \(e-1\) ticks are affine, followed by a division from
\((h,\ell)=(e,e)\) to \((1,e-1)\).
Since \(e-1\ge2\), the subsequent instruction word repeats
\(e-1\) affine ticks and one division. Counting these ticks gives the formula.
\end{proof}

\begin{theorem}[Complete homogeneous classification]
For \(q=2,c=0,d=1+\sqrt{-D}\):
\begin{enumerate}
\item if \(D=1\) or \(D=3\), every nonzero seed is eventually periodic
with primitive period twelve;
\item for every other \(D\ge1\), every nonzero seed escapes along the
full raw sequence.
\end{enumerate}
\end{theorem}
\begin{proof}
A legal homogeneous word with \(A\) affine instructions and \(B\)
divisions taking \(x\) to \(y\) satisfies \(d^B y=2^A x\), by induction
on the actual guarded word. A nonempty return at a nonzero state
would therefore give \(d^B=2^A\). Neither count can be zero, since
norms would force the other count zero. For \(D\notin\{1,3\}\) the
preceding lemma excludes every return and every collision.
A disk contains only finitely many points of \(R_D\), because
\(|x|\le R\) and \(|y|\le R/\sqrt D\). The injective orbit visits
each disk only finitely often, proving full-sequence escape.
The preceding quantitative theorem gives a stronger conclusion for
the dyadic cases \(e\ge3\); the norm-prime budget gives an eventual
doubling tail whenever \(D+1\) has an odd prime factor.

Repeated exact divisions terminate for every nonzero seed, since
each decreases the positive integral norm by the factor \(M>1\).
Let \(t\) be the resulting \(d\)-free state.

For \(D=1\), \(2=-i d^2\). Each affine tick raises the \(d\)-power
from zero to two and multiplies the phase by \(-i\); the next two
ticks divide. The three norm levels are distinct, and the phase
has exact order four. Thus the resulting cycle has primitive period
\(3\cdot4=12\).

For \(D=3\), write \(t=x+yw\). If \(x-y\) is odd, the following list
is one legal twelve-tick cycle, with the final point equal to the first:
\[
\begin{gathered}
t,\ 2t,\ 4t,\ (1-w)t,\ 2(1-w)t,\ -(1+w)t,\\
-t,\ -2t,\ -4t,\ -(1-w)t,\ -2(1-w)t,\ (1+w)t,\ t.
\end{gathered}
\]
The six division inputs are
\(4t,2(1-w)t,-(1+w)t,-4t,-2(1-w)t,(1+w)t\).
Their coordinate differences are, respectively,
\(4(x-y),4(x+y),4y,-4(x-y),-4(x+y),-4y\).
The other inputs have differences
\(x-y,2(x-y),2(x+y),-(x-y),-2(x-y),-2(x+y)\),
none divisible by four since \(x-y\) and \(x+y\) are odd.
The displayed transitions now follow from \(d(1-w)=4\).

If \(x-y\equiv2\pmod4\) and \(x,y\) are even, put \(u=t/2\).
Then \(u\) has odd coordinate difference and \(t=2u\) lies on its cycle.
If \(x,y\) are odd, put \(u=-dt/4\).
Both coordinates of \(dt\) are divisible by four, and the coordinate
difference of \(u\) is \(y\), which is odd. Since \(d\bar d=4\),
\(t=-\bar d u\), again a point of the cycle.
These cases exhaust all \(d\)-free states.

Finally any nonempty return for \(D=3\) has \(A=B\) by norms, and
\((d/2)^B=1\). The element \(d/2\) has exact order six
(\((d/2)^3=-1\)). Thus \(B\) is a positive multiple of six and the
return time \(A+B\) is a positive multiple of twelve.
The constructed twelve-tick return is primitive.
\end{proof}

\begin{corollary}[Nondyadic divisions]
Suppose \(D+1\) is not a power of two. Every nonzero homogeneous seed
has only finitely many divisions and then doubles at every raw tick.
\end{corollary}
\begin{proof}
The odd part of \(D+1\) exceeds one. A least divisor greater than one
is a prime \(p\); otherwise its proper factors contradict minimality.
Thus \(p\mid D+1=\OrderNorm(d)\) and \(p\nmid4=\OrderNorm(q)\).
The norm-prime budget applies to every nonzero seed.
\end{proof}

\section{Dyadic offsets and occupation}

\begin{theorem}[Odd-offset dyadic classification]
Let \(D=2^e-1\), \(e\ge1\), \(q=2\), \(d=1+\sqrt{-D}\), and
\(c=c_x+c_yw\in R_D\) with \(c_x-c_y\) odd.
The zero seed is fixed. A nonzero seed is eventually fixed if and
only if it belongs to
\[
\{-d^r c:r\in\N\}.
\]
Every such seed reaches the fixed point \(-c\).
Every other seed escapes along the full raw sequence.
\end{theorem}
\begin{proof}
A finite division prefix ends at a \(d\)-free point \(t\), with
\(z=d^r t\). After one affine tick, the coordinate difference is
\(2(t_x-t_y)+(c_x-c_y)\), which is odd. Every subsequent affine
difference is odd, so no division can occur. The initial point
is already \(d\)-free. Consequently, for every \(j\ge0\),
\[
\RawMap^j(t)=2^j(t+c)-c.
\]
If \(t=-c\), this is fixed; \(d\nmid c\) follows from its odd
coordinate difference. Otherwise
\[
|\RawMap^j(t)|\ge2^j|t+c|-|c|\longrightarrow\infty.
\]
The division-prefix identity proves both directions of the exact
exceptional-seed description. No independent branch law is used.
\end{proof}

\begin{corollary}[A division-frequency restriction]
For \(D=2^e-1\), \(e\ge3\), \(c=0\), and any nonzero seed, let
\(B_j\) count divisions in the first \(j\) ticks. Then
\[
B_j\ge\left\lfloor\frac{j}{e+1}\right\rfloor,\qquad
2j-(e+2)B_j\longrightarrow\infty.
\]
In particular,
\[
\frac1{e+1}\le\liminf_{j\to\infty}\frac{B_j}{j}
\le\limsup_{j\to\infty}\frac{B_j}{j}=\frac1e.
\]
\end{corollary}
\begin{proof}
Modulo \(d\), each affine tick doubles the coordinate difference.
After \(e\) such ticks it is zero modulo \(2^e\). Therefore every
block of \(e+1\) raw ticks contains a division.
Norm multiplicativity gives
\[
\OrderNorm(\RawMap^j(z))
=2^{\,2j-(e+2)B_j}\OrderNorm(z).
\]
Since \(D\notin\{1,3\}\), this norm tends to infinity by the homogeneous
classification. The exponent must tend to infinity. The lower count bound gives the first inequality. The dyadic valuation
theorem proves the stronger exact frequency limit \(1/e\).
\end{proof}

\begin{remark}
For \(D=7\), the norm relation \(4^3=8^2\) does not yield an element
return: \(d^2=-6+2\sqrt{-7}\ne8\). Every nonzero trajectory escapes,
yet divisions occur in every block of four ticks.
Thus a scalar norm balance alone does not classify critical cycles.
\end{remark}

\section{Exact finite figures}

The bounded arithmetic certificate checks \(372704\) entries.
The coordinate grid uses five orders, divisor coefficients in
\([-4,4]\), and seed coefficients in \([-6,6]\).
First-division checks use \(1\le D\le64\) and \(-256\le n\le256\).
The dyadic valuation checks cover \(672\) nonzero seeds for
\(e=3,4,5,6\). Independent integer-matrix references check
\(128\) dyadic instructions per seed on the smaller grid \([-4,4]\).
The infinite conclusions follow from the proofs, not these tests.

\begin{figure}[tbp]
\centering
\includegraphics[width=\linewidth]{../other/DyadicWordCounts.pdf}
\caption{Exact division counts for the seed one and
\(D=7,15,31,63\), through \(128\) raw ticks. For \(j\ge1\) these
are \(\lfloor(j-1)/e\rfloor\), not sampled frequencies from an
independent branch law. Each instruction word eventually has period
\(e\), while each arithmetic trajectory has positive escape exponent.}
\end{figure}

\begin{figure}[tbp]
\centering
\includegraphics[width=\linewidth]{../other/FirstDivisionDensity.pdf}
\caption{The proved first-division density \(1/\OddPart(D+1)\), sampled at
integer \(1\le D\le64\). Density one occurs at
\(D=1,3,7,15,31,63\), although only \(D=1,3\) have homogeneous cycles.
Connecting segments represent discrete samples, not a continuous
parameter law.}
\end{figure}

\begin{figure}[tbp]
\centering
\includegraphics[width=\linewidth]{../other/FirstDivisionSurface.pdf}
\caption{The exact integer-grid first-division mesh for
\(1\le D,n\le12\), projected by
\(\xi=D+n/2\), \(\eta=n/2+2h\).
Here \(h=a+1\) for a first divisible state at index \(a\), and \(h=0\)
when division never occurs. Projection can hide or superimpose points;
no continuous surface theorem is inferred.}
\end{figure}

\begin{figure}[tbp]
\centering
\includegraphics[width=\linewidth]{../other/EqualNormDynamics.pdf}
\caption{Gaussian homogeneous trajectories with \(q=2+i\), starting at
one. Taking \(d=q\) gives the exact two-cycle \(1,q,1\).
Taking \(d=2-i\) gives a division-free trajectory: modulo \(d\),
\(i=2\) and \(q=4\) in \(\Z/5\Z\), so powers of \(q\) never vanish.
In both cases \(\OrderNorm(q)=\OrderNorm(d)=5\).
The finite height traces display an exact guard distinction that norms
alone cannot recover.}
\end{figure}

\clearpage
\bibliographystyle{amsalpha}
\bibliography{QuadraticArithmetic}
\end{document}
